\chapter[State of the Art]{State of the Art and Technical Background}

This chapter keeps the state of the art in the same role it has in the paper:
it defines the nearest research families and the fair-comparison caveats, rather
than providing an encyclopaedic survey of every adjacent radio-map method. The
longer survey draft remains in \texttt{state\_of\_art.tex} as source material,
but the compiled thesis now keeps only the material needed to justify the final
prior-anchored CKM model.

\section{Channel knowledge maps and dense radio-map prediction}
\label{sec:soa_ckm_context}

Channel knowledge maps (CKMs) store location-indexed radio information that can
be queried without running a full ray tracer or a new measurement campaign for
every candidate link \cite{zengx2021ckm}. Dense path-loss maps are the most
common target, but the same idea extends naturally to delay spread and angular
spread when the map is used to reason about wideband or spatial channel quality.
The thesis problem is therefore closer to dense radio-map generation than to a
single-link path-loss regression task: one input topology and one UAV antenna
height must produce a full \(513\times513\) map.

Classical propagation models supply the physical vocabulary used by the final
method. Free-space path loss and two-ray propagation describe the LoS radial
structure; COST231/Hata-like and A2G models provide empirical NLoS attenuation
shapes; 3GPP and WINNER-style models motivate log-domain treatment of spread
statistics and their dependence on visibility, range, elevation and local
clutter \cite{rappaport,goldsmith,cost231,alhourani2014,tr38901,winner2}.
These models are not accurate enough as fixed formulas for the CKM dataset, but
they are useful as calibrated priors once the coefficients are fitted only on
training cities.

\section{Deep learning, hybrid models and distribution-aware heads}
\label{sec:soa_dl}

Deep radio-map prediction is usually formulated as image-to-image regression.
RadioUNet \cite{radiounet2020}, PMNet \cite{pmnet2023}, RadioGUNet
\cite{radiogunet2025}, RMTransformer \cite{rmtransformer2025}, PathFinder
\cite{pathfinder2025}, Gao et al.\ \cite{gao2026}, ReVeal \cite{reveal2025},
and AIRMap \cite{airmap2025} are the closest methodological neighbours. They
show that CNN, graph, attention, corridor-weighting, and digital-twin
approaches can produce accurate radio maps when the geometry and split
assumptions match the task. However, most of these references focus on path
loss or received power only, often with fixed or terrestrial transmitter
settings, and do not simultaneously predict dense path loss, delay spread and
angular spread for a centred but variable-height UAV.

Physics-informed and hybrid models are relevant because pure image regression
can waste capacity relearning radial propagation laws. Prior-residual learning
instead writes the prediction as
\[
  \hat y(x) = p(x) + \Delta_\theta(x),
\]
where \(p(x)\) is a frozen physical or calibrated prior and
\(\Delta_\theta\) is a learned residual. This idea only helps when the prior is
strong and auditable; otherwise the network inherits a poor bias. The final
method therefore calibrates the priors on training cities, freezes them, and
then asks the neural model to improve the residual.

Distribution-aware prediction is the other essential ingredient. A single MSE
head assumes that the conditional target is well summarized by its mean. The
late project diagnostics showed that this was false: LoS and NLoS path-loss
marginals are different, and delay/angular spread maps contain a near-zero LoS
spike plus NLoS-heavy tails. Heteroscedastic losses, mixture-density networks,
GMM heads and diffusion-style generative models all address this mismatch in
different ways \cite{kendall2017uncertainties,saleh2021probabilistic,
garciamarti2020mixture,radiodiff2025}. The final model keeps the practical
part of that lesson: a prior-anchored residual GMM head emits mixture
parameters around the frozen prior, then the selected RMSE-dominant checkpoint
exports a local-weighted point map rather than making a calibrated
probabilistic-density claim.

\section{Fair comparison protocol}
\label{sec:soa_benchmarks}

Headline RMSE values across radio-map papers are not directly interchangeable.
The reported number changes with target definition, map size, frequency,
transmitter setting, split discipline, whether building pixels are counted, and
whether the task is dense map generation or sparse reconstruction. Table
\ref{tab:fairness} lists the comparison axes that matter for this thesis.

\begin{table}[htbp]
\centering
\caption{Axes required for a fair reading of radio-map metrics.}
\label{tab:fairness}
\scriptsize
\setlength{\tabcolsep}{3pt}
\begin{tabularx}{\textwidth}{p{0.24\textwidth}X}
\toprule
Axis & Why it matters here \\
\midrule
Target and units & PL in dB, DS in ns and AS in degrees have different scales and target distributions. \\
Dense vs sparse input & Dense generation from topology is harder to compare with sparse-measurement interpolation. \\
Split policy & City-holdout is stricter than random sample or random pixel splits because it tests new urban layouts. \\
Receiver mask & Building pixels are not receiver locations; including them can make metrics artificially easier. \\
Transmitter setting & The thesis uses a horizontally centred UAV with variable antenna height, not arbitrary terrestrial transmitters. \\
Height handling & UAV height changes distance, reflection, elevation and local clearance; it is a core variable, not metadata. \\
Runtime path & A surrogate must be compared against a clearly specified ray-tracing path and hardware. \\
\bottomrule
\end{tabularx}
\end{table}

\section{Position of the final model}
\label{sec:soa_position}

The closest state-of-the-art references are useful as scale checks, not as a
single leaderboard. The final result is strongest when stated narrowly: under
the CKM city-holdout protocol, a prior-anchored residual GMM model improves
strong frozen priors for dense PL/DS/AS maps, while keeping a fast generator
path. Table~\ref{tab:position} summarizes that position.

\begin{table}[htbp]
\centering
\caption{Compact state-of-the-art positioning.}
\label{tab:position}
\scriptsize
\setlength{\tabcolsep}{3pt}
\begin{tabularx}{\textwidth}{p{0.23\textwidth}p{0.25\textwidth}X}
\toprule
Family / reference & Relevant metric or capability & Caveat for this thesis \\
\midrule
RadioUNet \cite{radiounet2020} & Dense radio-map CNN, about \SI{1.6}{dB}--\SI{3.1}{dB} PL-scale error in its setting & No variable-height UAV and no DS/AS targets. \\
PMNet / ICASSP \cite{pmnet2023,icassp2023challenge} & Strong outdoor path-loss-map baseline, about \SI{1.38}{dB} after score conversion & Different benchmark and no joint spread prediction. \\
RadioGUNet \cite{radiogunet2025} & Graph/U-Net radio maps, \SI{1.304}{dB} DPM and \SI{1.936}{dB} IRT reported & Outdoor RadioMapSeer setting; not the CKM UAV protocol. \\
ReVeal \cite{reveal2025} & Outdoor sparse-measurement mapping, \SI{1.95}{dB} RMSE & Sparse reconstruction rather than dense topology-only generation. \\
RMTransformer / PathFinder \cite{rmtransformer2025,pathfinder2025} & Attention/domain-adaptation radio-map references around the same PL scale & Different split and transmitter assumptions. \\
Gao et al.\ \cite{gao2026} & Outdoor corridor-weighted PL model, \SI{5.59}{dB} on ICASSP~2023 and \SI{12.19}{dB} on ITU~2024 & Recent comparable PL reference; final PL is numerically lower, but protocols differ. \\
AIRMap \cite{airmap2025} & Digital-twin radio-map generation with path-gain error below \SI{4}{dB} & Different target, map size and evaluation path. \\
\textbf{Final model (this thesis)} & \textbf{\SI{1.6519}{dB} PL, \SI{26.5570}{ns} DS, \(11.3854^\circ\) AS} & \textbf{Strict city holdout, \(513\times513\) dense maps, centred variable-height UAV and joint PL/DS/AS output.} \\
\bottomrule
\end{tabularx}
\end{table}

The novelty is therefore not that a CNN is used for radio maps. The novelty is
the combination of train-only height-aware priors, dense UAV CKM targets,
LoS/NLoS and morphology regime handling, and a residual GMM head that improves
rather than replaces the frozen estimators.
